\scp{Summary of issues and implications}{08-05}
\tsc{Taliabu} historical phonology (\srf{08-01}),
morphology (\srf{08-02}), and vocabulary (\srf{08-03}) are significantly
different from \tsc{Sula-Buru}, and in particular from Buru to a degree that is far different
from what one expects for “closely related” languages,
and below the range of lexical similarity for languages that
are normally included in the same subgroup.

Given that \tsc{Taliabu} appears to have links both to \tsc{Celebic}
languages to the west
and \tsc{Sula-Buru} languages to the southeast,
we must ask what is more basic,
and what is more easily explained by diffusion
and contact? The consistency
and regularity of *-ay > \it{e},
* aw > \it{o},
*j > \it{y}, shared with \tsc{Celebic} languages,
and *ə > \it{o} shared with \tsc{Eastern Celebic} languages,
and not shared with any other \tsc{Sula-Buru},
\tsc{Ambon-Seram} or \tsc{Seram-Tanimbar-Bomberai} languages
is striking and hard to account for by contact or diffusion.
While a broad pattern of *R > \it{h} is shared with \tsc{Sula-Buru}
languages (although Mangole-Sanana *R > {\0}),
in the totality of the available data it is not shared in the
same way in the same words.
For example, in \trf{08-11} all \tsc{Taliabu}
and \tsc{Sula-Buru} languages have their expected sound correspondences
for *R for *uRat ‘vein’,
*baqəRu ‘new’, *Raya ‘big’,
*qabaRa ‘shoulder’, *ñaRa ‘brother (f.s.)’, and *tabuRi ‘conch’.
But for *maRuqanay ‘male’,
\tsc{Taliabu} languages do not have their expected *R > \it{h}.
Similarly for *Rusuk ‘rib’,
Buru has the expected \it{hoso-n} ‘ribcage’,
whereas Soboyo has unexpected \it{usuk} ‘rib’.

Looking at things bottom-up from the perspective of the surface
forms and asking “where do these languages get \it{h}?” we find great variety.
Buru-Lisela only have  *R > \it{h}.
Sanana-Mangole have *R, *s > \it{h}.
Sanana also has *d > \it{h}.
Ambelau has *R, *l, *d > \it{h}.
\citet{Blust1981} observes that Soboyo \it{h} is from multiple
sources, such as PMP *h,
*p, *R, *d, *r, *b.
Buru consistently preserves the distinctions of those proto phonemes.

\tsc{Saluan-Banggai} merges *j/*y > **y.
\tsc{Taliabu} has a partial merger of *z/*j/*y > **y with lexically specific exceptions.
Thus they both preserve *y = **y (\trf{08-08}).
\tsc{Sula-Buru} also preserves *y = **y,
with Buru-Lisela **y > {\0} as a secondary development.
In \tsc{Sula-Buru} *j and *z do not shift to **y.
Retentions are not useful for subgrouping,
so it is the merger of *j/*y > **y that links \tsc{Taliabu} with
\tsc{Saluan-Banggai}, and distinguish it from \tsc{Sula-Buru}.

Putting it all together looking at significant differences in
the historical phonology, different morphology, significant differences in the lexicon
in both form and frequency, and \citeauthor{Bloyd2015}'s (\citeyear{Bloyd2015,Bloyd2020})
avoidance of trying to include
\tsc{Taliabu} data with his Proto-Sula,
there seem to be a bundle of issues making a strong case for
saying the \tsc{Taliabu} node has a very different prehistory
than \tsc{Sula-Buru} and should not be included with other languages of ‘Central Maluku’.
Other than PMP *R > \it{h},
the consistency and regularity of the bundle of exclusive innovations
in the historical phonology shared between \tsc{Taliabu}
and \tsc{Eastern Celebic} seems quite different to the more haphazard
and often lexically specific nature of what is shared with \tsc{Sula-Buru}.
\citet{Mead2003a} identifies the following innovations defining
\tsc{Celebic}. 

\begin{enumerate}[noitemsep]
		\item	*C\sub{1}C\sub{2} > C\sub{2} (where C\sub{1} is not a nasal)
		\item	*h > {\0}
		\item	*d > **r
		\item	*-ay > \it{e}
		\item	*-aw > \it{o}
		\item	*j > **y, {\0}
\end{enumerate}

Of these, {\#}1 has also been claimed to define all of “CEMP”
\citep[246]{Blust1993}, although we discuss problems with that claim in \srf{13-02-01},
including its presence defining \tsc{Celebic} languages.
\tsc{Taliabu} shares {\#}4-6 (see \srf{08-01} above,
with {\#}6 having a couple of lexically specific exceptions).
Regarding {\#}3, we have shown that there is a path by which
that is also shared: *d > **r > \it{h}.
While Soboyo might show some cases of PMP *h = \it{h} \citep{Blust1981},
the normal reflex for Soboyo
and \tsc{Taliabu} is PMP *h > {\0} (see \srf{13-02-02} for discussion
of the sporadic areal phenomenon of irregular C-onsets on V-initial
roots, with {\#}h, {\#}w, {\#}y, {\#}l, {\#}t, {\#}s, etc).

\citet{Mead2003a,Mead2003b} identifies the following innovations defining
\tsc{Eastern Celebic} languages: 

\begin{enumerate}[noitemsep]
	\item	*ə > \it{o}
	\item	*-iq > **eq
	\item	antepenultimate *a > *\it{o}
\end{enumerate}

Of these, in \srf{08-01} we showed that *ə > \it{o} in \tsc{Taliabu}.
There is not enough data to say whether {\#}2
and {\#}3 are found in \tsc{Taliabu}.
\citet{Mead2003b} identified the following innovations as defining
\tsc{Saluan-Banggai}.

\begin{enumerate}[noitemsep]
	\item	*awa > **oa
	\item	*b{\#},*g{\#} > **p{\#}, **k{\#} (devoicing)
	\item	*q > \it{ʔ} 
\end{enumerate}

Of these, in \srf{08-01} we showed that there is one instance
of *awa > **oa in \tsc{Taliabu}.
With limited data, it appears that *b{\#},*g{\#} > **p{\#},
**k{\#} > {\0}, is consistent with the loss of many final consonants.
Similarly for *q > **ʔ > {\0}.

So we tentatively group \tsc{Taliabu} with \tsc{Celebic},
perhaps in its own node within \tsc{Eastern Celebic} node,
or as a sister node to \tsc{Eastern Celebic},
noting that it shares some diagnostic sound changes with \tsc{Saluan-Banggai},
but with the caveat that more data
and more research are clearly required.

The presence of the preposed possessor in Banggai associated
with OV typology (not found in other \tsc{Saluan-Banggai} languages)
suggests either a pre-Austronesian SOV (Papuan) substrate,
or else significant contact with other Austronesian languages
that were already impacted by such a Papuan substrate (see \srf{02-03-01-01}).

Furthermore, the shared phonological prehistory that links the
\tsc{Taliabu} node with \tsc{Celebic} languages,
along with different morphology
and significant differences in the lexicon that also set it apart
from all other \tsc{Sula-Buru}
and \tsc{Ambon-Seram} languages, requires that some higher level
claims about Austronesian dispersal be re-examined.
In \srf{02-01} it was noted that \citet{Blust2012b} rejects \citeauthor{Schapper2011}’s
\citeyear{Schapper2011} “implicit claim” \citep[273]{Blust2012b} of how the Austronesians
encountered the cuscus and bandicoot, because it
“implies an entry via Sulawesi into
the Lesser Sundas, with subsequent eastward movement into the Moluccas.
If this were the case,
we would expect some linguistic trace in Sulawesi,
but there is none; [{\ldots}] about possible subgrouping ties
between at least some of the languages of Sulawesi
and languages in eastern Indonesia,
closer scrutiny of the evidence reveals nothing of substance
to support this idea” \citep[274]{Blust2012b}.

Yet scrutiny of the \tsc{Taliabu} data from the historical phonology
do show that there was at least one incursion into Maluku from
neighbouring islands in Sulawesi.
In addition, the lexical
and morphological data show that this incursion from Sulawesi
had significant subsequent impact into the northern part of the
\tsc{Sula-Buru} subgroup, particularly in Mangole and Sanana, but also reaching Lisela,
Buru and Ambelau to lesser degrees.

